\documentclass[zihao=-4, a4paper, oneside]{ctexbook}
\usepackage[margin=2.8cm]{geometry}
\usepackage{amsmath, amssymb}
\usepackage{booktabs}
\usepackage{graphicx}
\usepackage{setspace}
\usepackage{titlesec}
\usepackage{hyperref}

\onehalfspacing

% ==== 学校信息（替换为自己的） ====
\newcommand{\university}{示例大学}
\newcommand{\thesisname}{本科毕业设计（论文）}
\newcommand{\thesistitle}{基于深度学习的示例研究}
\newcommand{\thesisauthor}{张三}
\newcommand{\thesisid}{2022010101}
\newcommand{\thesisadvisor}{李四 教授}
\newcommand{\thesismajor}{计算机科学与技术}
\newcommand{\thesisdate}{2026年6月}

\begin{document}

% ================= 封面 =================
\begin{titlepage}
  \centering
  \vspace*{2cm}
  {\zihao{1}\heiti \university}\\[1cm]
  {\zihao{2}\heiti \thesisname}\\[3cm]
  \zihao{3}
  \begin{tabular}{rl}
    \textbf{题\hspace{2em}目：} & \thesistitle \\[0.5em]
    \textbf{姓\hspace{2em}名：} & \thesisauthor \\[0.5em]
    \textbf{学\hspace{2em}号：} & \thesisid \\[0.5em]
    \textbf{专业：} & \thesismajor \\[0.5em]
    \textbf{指导教师：} & \thesisadvisor \\
  \end{tabular}
  \vfill
  {\zihao{3} \thesisdate}
\end{titlepage}

% ================= 摘要 =================
\frontmatter
\chapter{摘\quad 要}
本文是一份中文毕业论文模板。摘要应概括研究背景、研究问题、
所提方法、实验结果与主要结论，中文摘要一般不少于 500 字。

\textbf{关键词：} 毕业论文；\LaTeX；模板

\chapter{Abstract}
This is the English abstract of the thesis. Replace it with a translation
of your Chinese abstract.

\textbf{Keywords:} thesis; LaTeX; template

% ================= 目录 =================
\tableofcontents

% ================= 正文 =================
\mainmatter

\chapter{绪论}
\section{研究背景与意义}
介绍选题背景、国内外研究现状以及研究的意义。
\section{研究内容}
概述本文的主要研究内容与组织结构。

\chapter{相关工作}
\section{方法 A}
综述相关研究工作~\cite{ref1}。
\section{方法 B}
补充其他相关工作~\cite{ref2}。

\chapter{方法}
\section{问题定义}
给出问题的形式化定义：
\begin{equation}
  \min_{\theta} \; \frac{1}{N} \sum_{i=1}^{N} \ell\!\left(f_\theta(x_i), y_i\right).
\end{equation}

\section{模型设计}
详细描述所提出的方法与模型结构。

\chapter{实验与分析}
\section{实验设置}
描述数据集、评价指标、对比方法与实现细节。

\begin{table}[htbp]
  \centering
  \caption{主要实验结果。}
  \begin{tabular}{lcc}
    \toprule
    方法 & 准确率（\%） & F1 \\
    \midrule
    基线 & 81.2 & 0.79 \\
    本文 & \textbf{88.1} & \textbf{0.87} \\
    \bottomrule
  \end{tabular}
\end{table}

\section{结果分析}
分析实验结果并讨论局限性。

\chapter{总结与展望}
总结全文工作，并展望未来研究方向。

% ================= 参考文献 =================
\backmatter
\begin{thebibliography}{9}
\bibitem{ref1}
作者一, 作者二. 相关论文标题[J]. 期刊名称, 2024, 12(3): 45--58.
\bibitem{ref2}
Author A, Author B. Another related paper[C]//Proceedings of a Conference.
2025: 100--110.
\end{thebibliography}

% ================= 附录 =================
\appendix
\chapter{补充材料}
附录内容：推导细节、补充图表、代码清单等。

\end{document}
